\begin{lemma}[Existence of subgradients of a convex function]
Let $E$ be a finite-dimensional real inner-product space, let $\Omega\subseteq E$
be open, let $\psi:E\to\R$ be convex on $\Omega$ (so in particular $\Omega$ is
convex) and let $z\in\Omega$. Then $\psi$ has a subgradient at $z$: there is
$q\in E$ with
\[
  \psi(z)+\langle q,x-z\rangle\le\psi(x)\qquad(x\in\Omega).
\]
\end{lemma}
\begin{proof}
\emph{Step 1: a ball on which $\psi$ is Lipschitz.}
Since $\Omega$ is open and $\psi$ is convex and real-valued on $\Omega$, $\psi$
is continuous on $\Omega$ and locally Lipschitz on $\Omega$ (preamble). Let
$U\subseteq\Omega$ be a neighbourhood of $z$ and $L\ge0$ a constant with
$|\psi(a)-\psi(b)|\le L|a-b|$ for $a,b\in U$. Choose $\rho>0$ so small that
$\overline B(z,2\rho)\subseteq U$; this is possible because $U$ is a
neighbourhood of $z$. Then $\psi$ is $L$-Lipschitz on $\overline
B(z,2\rho)\subseteq\Omega$.

\emph{Step 2: differentiability points accumulating at $z$.}
The ball $B(z,2\rho)$ is open and $\psi$ is Lipschitz on it, so by Rademacher's
theorem (preamble) $\psi$ is differentiable at almost every point of
$B(z,2\rho)$; write $D$ for the set of those points of $B(z,2\rho)$ at which
$\psi$ is differentiable, so that $B(z,2\rho)\setminus D$ is null. Fix an
additive Haar measure $\mu$ on $E$. For each integer $j\ge1$ the ball
$B(z,\rho/j)$ has $\mu(B(z,\rho/j))>0$ (preamble), so it cannot be contained in
the null set $B(z,2\rho)\setminus D$; hence there is
$w_j\in D$ with $|w_j-z|<\rho/j$. In particular $w_j\to z$ and
$w_j\in B(z,\rho)$ for every $j$.

\emph{Step 3: at a differentiability point the gradient is a subgradient.}
Let $w\in D$ and let $x\in\Omega$. For $t\in(0,1]$ the point $w+t(x-w)$ lies in
$\Omega$, since $\Omega$ is convex and $w,x\in\Omega$, and convexity of $\psi$ on
$\Omega$ with the weights $1-t$ and $t$ gives
\[
  \psi\bigl(w+t(x-w)\bigr)\le(1-t)\psi(w)+t\psi(x),
\]
that is,
\[
  \frac{\psi(w+t(x-w))-\psi(w)}{t}\le\psi(x)-\psi(w).
\]
As $t\downarrow0$ the left-hand side tends to $\langle\nabla\psi(w),x-w\rangle$,
because $\psi$ is differentiable at $w$ with gradient $\nabla\psi(w)$ and
$\psi(w+th)=\psi(w)+t\langle\nabla\psi(w),h\rangle+o(t)$ for fixed $h=x-w$.
Therefore
\begin{equation}\label{eq:cse-subgrad-at-w}
  \psi(w)+\langle\nabla\psi(w),x-w\rangle\le\psi(x)\qquad(w\in D,\ x\in\Omega).
\end{equation}

\emph{Step 4: the gradients at the $w_j$ are bounded by $L$.}
Fix $j$ and let $e\in E$ with $|e|=1$. For $0<t\le\rho$ the point $w_j+te$
satisfies $|w_j+te-z|\le|w_j-z|+t<\rho+\rho=2\rho$, so $w_j+te\in
B(z,2\rho)\subseteq\Omega$. Applying \eqref{eq:cse-subgrad-at-w} with
$x=w_j+te$, and then the Lipschitz bound of Step 1,
\[
  t\langle\nabla\psi(w_j),e\rangle\le\psi(w_j+te)-\psi(w_j)\le Lt .
\]
Dividing by $t>0$ gives $\langle\nabla\psi(w_j),e\rangle\le L$ for every unit
vector $e$. Taking $e=\nabla\psi(w_j)/|\nabla\psi(w_j)|$ when
$\nabla\psi(w_j)\ne0$ yields $|\nabla\psi(w_j)|\le L$; and the bound is trivial
when $\nabla\psi(w_j)=0$. So $|\nabla\psi(w_j)|\le L$ for every $j$.

\emph{Step 5: passing to the limit.}
The sequence $(\nabla\psi(w_j))_j$ is bounded, so by Bolzano--Weierstrass
(preamble) it has a subsequence $(\nabla\psi(w_{j_i}))_i$ converging to some
$q\in E$. Let $x\in\Omega$. By \eqref{eq:cse-subgrad-at-w},
\[
  \psi(w_{j_i})+\langle\nabla\psi(w_{j_i}),x-w_{j_i}\rangle\le\psi(x)
  \qquad(i\ge1).
\]
As $i\to\infty$ we have $w_{j_i}\to z$, hence $\psi(w_{j_i})\to\psi(z)$ by the
continuity of $\psi$ on $\Omega$ (Step 1), and
$\langle\nabla\psi(w_{j_i}),x-w_{j_i}\rangle\to\langle q,x-z\rangle$ by the
joint continuity of the inner product. Therefore
$\psi(z)+\langle q,x-z\rangle\le\psi(x)$. Since $x\in\Omega$ was arbitrary, $q$
is the required subgradient.
\end{proof}
